% Source pins (SHA-256): PifoStatement.lean=fe9475ebfa2c462d27cb0a99781926074e1281396f2c9e8d32d6dfca788d814e; pifo-normalization.md=2d287de7924aa30479218da3e0a9a31577ecb97c5a34395c089f7abed98f61e7.
\section{Normalization and canonical forms}
\label{sec:normalization-canonical}

The proof of \autoref{thm:main} joins two schedulers through a decomposition
chosen from both of them.  A stronger consequence is that each scheduler has
its own canonical decomposition tree.  The tree is finite, is recoverable from
flush behavior, and completely characterizes interleaved behavior.

Throughout this section \(A=\Fin(k)\), with its fixed enumeration, and all
schedulers are valid.  For \(X\subseteq A\), write \(S|_X\) for the scheduler
obtained by retaining only the types in \(X\).  Runs whose pushes lie in \(X\)
are unchanged by this restriction.

\subsection{The local data}

For distinct \(x,y\in A\), let
\[
  \mathsf{eff}_S(x,y)\in\{<,=,>\}
\]
be the effective comparison of their two-type restriction.  Concretely,
follow their assigned paths together through every common child.  At their
first different child, compare the two reference ranks; if they instead end
at one leaf, compare their leaf ranks.  By \lemref{lem:two-type}, this one
comparison determines the restricted scheduler on every operation word.  It
is local under restriction and is recovered from the two drains in
\autoref{tab:binary-recovery}.

Triples carry one further bit.  For comparisons \(c,d\in\{<,=,>\}\), let
\(\operatorname{comp}(c,d)\) be their transitive composition: composing with
\(=\) returns the other comparison, two \(<\)'s compose to \(<\), two
\(>\)'s compose to \(>\), and opposite strict signs give \(\None\).

\begin{definition}[Rooted-triple obstruction]
\label{def:rooted-obstruction}
For pairwise-distinct \(x,y,z\), restrict \(S\) to these three types and use
the triple-flattening lemma to replace that restriction by an
interleaved-equivalent flat scheduler.  Write \(\ell(u)\) for the leaf of
type \(u\) in this representative.  Then \(R_S(x,y\mid z)\) holds when
\begin{equation}
\begin{split}
  &\ell(x)=\ell(y)\ne\ell(z),\\
  &\operatorname{comp}
      \bigl(\mathsf{eff}_S(x,z),\mathsf{eff}_S(z,y)\bigr)=\Some(c),
    \qquad c\ne\mathsf{eff}_S(x,y).
\end{split}
\label{eq:rooted-obstruction}
\end{equation}
The triple lemma makes this predicate independent of the chosen flat
representative.  It is also local under every further restriction containing
\(x,y,z\).
\end{definition}

The structural construction of the required flat representative is given in
\autoref{sec:structural-decision}; its \(2+1\) case is important because the two
types in the shared branch must be collapsed to one synthetic leaf.

For \(X\subseteq A\), form an undirected graph \(G_S(X)\) on \(X\) by putting
an edge between distinct \(x,y\) exactly when some
\(z\in X\setminus\{x,y\}\) satisfies
\begin{equation}
  R_S(x,y\mid z)\ \lor\ R_S(y,x\mid z).
  \label{eq:obstruction-edge}
\end{equation}
Let \(\Pi_S(X)\) be the partition of \(X\) into connected components of this
graph.

\begin{lemma}[Locality and properness]
\label{lem:canonical-properness}
The data \(\mathsf{eff}_S\) and \(R_S\) on \(X\) agree with the corresponding
data of \(S|_X\).  If \(|X|>1\), every block of \(\Pi_S(X)\) is a proper
subset of \(X\).
\end{lemma}

\begin{proof}
Pair locality follows by suppressing path prefixes used by both types; triple
locality follows from the same common-prefix suppression and the
triple-flattening lemma.  For properness, descend through every node at which
all types in \(X\) use one child.  If this reaches a leaf, every triple lies in
one flat leaf and \autoref{eq:rooted-obstruction} is false, so the graph has no
edges.  Otherwise the first branching node partitions \(X\) into at least two
nonempty child alphabets, and no obstruction edge crosses those alphabets.
Thus \(G_S(X)\) is disconnected in either case.
\end{proof}

The component partition has exactly the compatibility needed at a canonical
root.

\begin{lemma}[Cell feasibility]
\label{lem:canonical-cell-feasibility}
For every \(X\subseteq A\) with \(|X|>1\), there is a rank function
\(q_X:X\to\Nat\) such that
\begin{equation}
  \cmp_{q_X}(x,y)=\mathsf{eff}_S(x,y)
  \qquad
  \text{whenever \(x,y\) lie in different \(\Pi_S(X)\)-blocks.}
  \label{eq:canonical-cell-feasibility}
\end{equation}
There is no constraint on the root comparison of two types in one block.
\end{lemma}

The lemma is the cell-feasibility part of the top-decomposition argument: one
normalizes each child order in turn, while the strong preservation clause
keeps all previously treated components fixed.  The resulting finite total
preorder is represented by natural ranks.  The absence of a within-block
constraint is essential---that comparison is resolved by the child and is
not an observable root invariant.

\subsection{The component-decomposition normal form}

\begin{definition}[Abstract canonical tree]
\label{def:abstract-canonical-tree}
Define \(\mathcal N_X(S)\) by strong recursion on \(|X|\).
\begin{itemize}
\item If \(X=\varnothing\), it is the empty tree, represented by a zero-child
      node.
\item If \(X=\{x\}\), it is one leaf carrying \(x\).
\item If \(|X|>1\), let
      \(\Pi_S(X)=\{C_1,\ldots,C_m\}\).  The root is labelled by this
      partition and by the table
      \begin{equation}
        \mathsf{eff}_S\bigm|_{
          \{(x,y):[x]_{\Pi_S(X)}\ne[y]_{\Pi_S(X)}\}},
        \label{eq:canonical-node-label}
      \end{equation}
      and its children are
      \(\mathcal N_{C_i}(S|_{C_i})\), one for each block.
\end{itemize}
Put \(\mathcal N(S)=\mathcal N_A(S)\).  The recursion terminates by
\lemref{lem:canonical-properness}.
\end{definition}

For \(k>0\), a concrete normal form \(N(S)\) is any scheduler realizing
\(\mathcal N(S)\): list siblings in a fixed order, use a witness from
\lemref{lem:canonical-cell-feasibility} at each internal node, and continue
through the recursively assigned child to the singleton leaf, where the
irrelevant leaf rank may be set to zero.  Such a scheduler is valid by
construction.

The empty alphabet needs a separate concrete convention.  Let
\begin{equation}
  E_0.\mathsf{topo}=\mathsf{node}[\,],
  \qquad
  E_0.\mathsf{assign}:\Fin(0)\to\mathsf{Path}
  \text{ the unique empty function},
  \qquad N(S)=E_0.
  \label{eq:empty-canonical-scheduler}
\end{equation}
This is valid vacuously.  Since an operation over \(\Fin(0)\) cannot push,
every operation is a pop and every pop from the initial state returns
\(\None\).  Hence all zero-type schedulers are interleaved-equivalent to
\(E_0\).

The concrete representative is canonical only up to inessential
presentation.  Equality of normal forms below means literal equality of the
labelled abstract trees \(\mathcal N\), or, equivalently, equality of concrete
realizations after identifying:
\begin{enumerate}
\item a permutation of sibling blocks together with the corresponding
      remapping of child indices;
\item any root-rank reparametrization preserving all cross-block comparisons,
      with within-block root ranks left free; and
\item any leaf-rank reparametrization preserving the induced weak order.
\end{enumerate}
Child indices do not break ties, so the first identification is operationally
invisible; queues inspect only ranks and arrival stamps.

\begin{theorem}[Canonical normalization]
\label{thm:canonical-normalization}
For every valid scheduler \(S\),
\begin{equation}
  S\intereq N(S).
  \label{eq:canonical-normalization}
\end{equation}
\end{theorem}

\begin{proof}
The \(k=0\) case was just proved.  If \(k=1\), unary active nodes merely
forward service, and the lone type is FIFO at its leaf; contracting those
nodes and setting the irrelevant leaf rank to zero gives the singleton leaf
\(N(S)\).  For \(k>1\), use strong induction on the support size and suppress
every unary active node.  If the reduced scheduler is a leaf, expand it to a
root ordered by the old leaf ranks and route each type to a singleton child;
this preserves every run.  Otherwise the canonical top-decomposition lemma
replaces the root by one child for each
\(C\in\Pi_S(A)\), uses a feasible root order satisfying
\autoref{eq:canonical-cell-feasibility}, and gives that child the behavior of
\(S|_C\).  Colored root congruence and stable child projection make this
replacement online-invisible.

Every \(C\) is proper, so the induction hypothesis replaces \(S|_C\) by
\(N(S|_C)\).  Interleaved equivalence is a congruence for the node
construction: the root couples only the component color, and paired children
then emit the same packets.  The resulting scheduler is \(N(S)\).
\end{proof}

\subsection{Flush determination and completeness}

\begin{lemma}[Flush determination]
\label{lem:canonical-flush-determination}
The pair \((\mathsf{eff}_S,R_S)\), and therefore \(\mathcal N(S)\), is
determined by the flush transformer \(F_S\).
\end{lemma}

\begin{proof}
For each distinct pair, the two drains of \(xy\) and \(yx\) recover the exact
comparison by \autoref{tab:binary-recovery}; both arrival orders are needed.  For
each distinct triple, restrict and relabel its types, flatten the restriction,
and use the finite three-label recovery lemma.  The drains of the injective
words of lengths two and three determine every rooted obstruction bit; the
singleton drains are fixed.  Thus all of \(\mathsf{eff}_S\) and \(R_S\) are
functions of \(F_S\).

At every recursive support \(X\), locality makes \(G_S(X)\), and hence
\(\Pi_S(X)\), a function of the restricted \(R_S\)-table.  The node label then
uses only the corresponding cross-block part of \(\mathsf{eff}_S\).  Strong
recursion therefore reconstructs the whole labelled tree.
\end{proof}

\begin{theorem}[Canonical characterization]
\label{thm:canonical-characterization}
For valid schedulers \(S,T\) over \(A\),
\begin{equation}
  S\intereq T
  \quad\Longleftrightarrow\quad
  \mathcal N(S)=\mathcal N(T)
  \quad\Longleftrightarrow\quad
  \bigl(\mathsf{eff}_S=\mathsf{eff}_T
        \ \land\ R_S=R_T\bigr).
  \label{eq:canonical-characterization}
\end{equation}
In particular, the canonical normal form is determined by flush behavior.
\end{theorem}

\begin{proof}
Equality of the two data tables forces equal canonical trees: locality gives
the same graph and component partition on every recursively reached support,
and equality of \(\mathsf{eff}\) gives the same cross-block label.  Conversely,
if the canonical trees agree, choose one concrete realization \(N_0\) of
their common tree.  By \autoref{thm:canonical-normalization},
\(S\intereq N_0\intereq T\).  This implies flush equivalence, so
\lemref{lem:canonical-flush-determination} recovers equal data tables.  This
semantic argument is important: an arbitrary triple need not itself occur as
a support in the canonical recursion, so \(R\) cannot in general be read
directly from one node of the tree.

Finally, interleaved equivalence implies flush equivalence and hence equal
data and equal trees.  The reverse implication is the common-realization
argument above.  These implications establish both equivalences in
\autoref{eq:canonical-characterization}.
\end{proof}
